% ============================================================
% Sección 6: RAG
% ============================================================
\section{RAG: Retrieval-Augmented Generation}

\begin{frame}{El problema que resuelve RAG}
  \begin{itemize}
    \item Un LLM solo ``sabe'' lo que vio en su entrenamiento: conocimiento
      \textbf{congelado en el tiempo} y sin datos privados/internos.
    \item Alternativas para darle información nueva:
      \begin{enumerate}
        \item Reentrenar/afinar el modelo (caro, lento).
        \item Meter todo en el prompt (no escala, límite de contexto).
        \item \textbf{RAG}: buscar solo lo relevante y ponerlo en el prompt.
      \end{enumerate}
  \end{itemize}
  \vfill
  \definicion{RAG}{
    Recuperar información relevante de una fuente externa \emph{en el
    momento de la consulta} y dársela al LLM como contexto para generar
    la respuesta.
  }
\end{frame}

\begin{frame}{Pipeline de RAG}
  \centering
  \begin{tikzpicture}[node distance=1.5cm, every node/.style={font=\small}]
    \node[cajaAgente] (q) {Pregunta\\del usuario};
    \node[cajaLLM, right=of q] (emb) {Modelo de\\embeddings};
    \node[cajaHerramienta, right=of emb] (db) {Base vectorial\\(documentos)};
    \node[cajaEntorno, below=of db] (ctx) {Fragmentos\\relevantes};
    \node[cajaLLM, below=of emb] (llm) {LLM};
    \node[cajaAccento, below=of q] (resp) {Respuesta\\final};

    \draw[flecha] (q) -- (emb);
    \draw[flecha] (emb) -- node[above, etiquetaFlecha]{vector consulta} (db);
    \draw[flecha] (db) -- node[right, etiquetaFlecha]{top-$k$ similares} (ctx);
    \draw[flecha] (ctx) -- (llm);
    \draw[flecha] (q) to[bend left=15] (llm);
    \draw[flecha] (llm) -- (resp);
  \end{tikzpicture}
  \vfill
  \footnotesize La pregunta viaja por dos caminos: se usa para \textbf{buscar}
  y se usa (junto al contexto recuperado) para \textbf{generar} la respuesta.
\end{frame}

\begin{frame}{¿Qué es un embedding, en una frase?}
  \definicion{Embedding}{
    Un vector de números que representa el \emph{significado} de un texto,
    de forma que textos con significado parecido quedan cerca en ese
    espacio vectorial.
  }
  \vfill
  \begin{itemize}
    \item Se indexan todos los documentos (o fragmentos/\emph{chunks}) como
      embeddings en una \textbf{base vectorial} (FAISS, Chroma, pgvector...).
    \item En consulta: se calcula el embedding de la pregunta y se busca por
      \textbf{similitud} (coseno) los fragmentos más cercanos.
  \end{itemize}
\end{frame}

\begin{frame}{Ejemplo 1: asistente sobre documentación interna}
  \small
  \begin{enumerate}
    \item Se indexan los manuales/wiki de la empresa en chunks de \textasciitilde500 tokens.
    \item Usuario pregunta: \emph{``¿Cuántos días de vacaciones tengo el primer año?''}
    \item Se recuperan los 3 fragmentos del manual de RR.HH. más relevantes.
    \item El LLM responde \textbf{citando} esos fragmentos, sin inventar la política.
  \end{enumerate}
  \vfill
  \begin{alertblock}{Por qué importa}
    Sin RAG, el modelo alucinaría una cifra genérica; con RAG, responde con
    la política \emph{real} de la empresa.
  \end{alertblock}
\end{frame}

\begin{frame}{Ejemplo 2: RAG dentro de un agente}
  \small
  \begin{enumerate}
    \item El agente recibe: \emph{``Resume los cambios de la última versión
      de nuestra librería''}.
    \item El agente decide llamar a la herramienta \texttt{buscar\_docs(query)}.
    \item La herramienta hace la recuperación (RAG) y devuelve fragmentos del changelog.
    \item El agente redacta el resumen usando esos fragmentos como evidencia.
  \end{enumerate}
  \vfill
  RAG aquí \textbf{no es un paso fijo}: es una herramienta más que el agente
  decide usar cuando la necesita.
\end{frame}

\begin{frame}{Retos comunes de RAG}
  \begin{itemize}
    \item \textbf{Chunking}: fragmentos muy grandes pierden precisión, muy
      chicos pierden contexto.
    \item \textbf{Recuperación pobre} = respuesta pobre, aunque el LLM sea excelente.
    \item Combinar búsqueda \textbf{semántica + por palabras clave} (\emph{hybrid search})
      suele dar mejores resultados.
    \item \emph{Re-ranking}: recuperar más candidatos y luego reordenar los
      mejores con un modelo más preciso (y más caro).
  \end{itemize}
\end{frame}
